\subsection{Cross-acquisition validation of a shared-recording clock}
\label{sec:p2-physionet-external-clock}
To determine whether the channel--clock mechanism is specific to one
motor-imagery acquisition, we reproduced the same frozen detector on
a distinct public EEG acquisition family: CHB--MIT scalp seizure EEG
(PhysioNet v1.0.0, dataset DOI 10.13026/C2K01R). The downloaded
digital EDF prefixes were directly decoded, using eight bipolar EEG
channels and 768 consecutive samples per source. The test sampled
public digital waveforms, \emph{not} independent clinically confirmed
duplicate episodes, and did not use seizure or clinical identity labels.

The source--target matcher required at least four distinct channels
with absolute Pearson correlation $\ge0.985$, a common temporal lag
within 100 samples and at least 230 overlapping samples; these numeric
settings were frozen in the previous EEGMMIDB experiment. Each
source generated eight known-reuse digital transformations (including
channel relabelling/permutation, common temporal shift, cropping,
channel omission, amplitude rescaling and low-amplitude noise) and
five controlled negative transformations. Among the negatives, an
\emph{asynchronous same-source} composite contained all eight source
channel contents but with eight individually chosen clock delays:
it was negative \emph{only} for the endpoint of a coherent
synchronised multichannel recording, and would be positive under
the different endpoint of any reused electrode samples.

\begin{table}[tb]
\centering\small
\caption{Source-frozen CHB--MIT test on real publicly released
digital EDF samples under \emph{experimentally imposed}
known-reuse transformations. The first two records were tested
before the subsequently untouched CHB03--CHB06 test group.
Multiple manipulations of one EDF do not constitute independent
clinical subjects. False matches are specific to the
coherent-multichannel-clock endpoint, not any-channel reuse.}
\label{tab:p2-chbmit-external-clock}
\setlength{\tabcolsep}{3pt}
\begin{tabular}{@{}p{0.40\linewidth}rrrr@{}}
\toprule
&\multicolumn{2}{c}{CHB01--02 (2 sources)}
&\multicolumn{2}{c}{CHB03--06 (new 4 sources)}\\
Method & TP/16 & FP/10 & TP/32 & FP/20\\
\midrule
Common clock + channel assignment &16&0&32&0\\
Independent best-lag channel Hungarian &16&2&32&4\\
Same-label zero-lag correlation &8&0&16&0\\
\bottomrule
\end{tabular}
\end{table}

The complete six-source outcome was 48/48 controlled positives and
zero of 30 constructed negative matches for the common-clock method,
compared with the same 48/48 positive detections but six of 30
negative matches for the unconstrained independent-lag Hungarian
comparator. Every discordance was an intentionally asynchronous
same-source example. This confirms the narrow practical value of
the \emph{global recording clock} constraint on another EEG task, not
new clinical classification accuracy, population-level false-positive
rates or algorithmic novelty over established multichannel
synchronisation systems. Source-level subgroup counts, original
HTTP byte-range status 206, prefix SHA-256 and independently
stored outcome records are available in
\href{https://github.com/lindicaphxag-tech/kaggle/actions/runs/38666039518}{CHB--MIT initial CI receipt}
and
\href{https://github.com/lindicaphxag-tech/kaggle/actions/runs/38666330505}{CHB--MIT four-new-source CI receipt}.
The full native EDF sources are distributed by PhysioNet and
are \emph{not} redistributed with this study.

Limitations include only six real digital source records from
one external acquisition family, author-constructed transformations,
heterogeneous semantics between channel reuse and common-clock
recording reuse, and a lack of baseline comparisons against
published multi-device synchronization, time-warping and
multichannel fingerprinting implementations.
